\begin{abstract}
Relações comerciais estão presentes no nosso cotidiano. Um vendedor expõe seus produtos à venda e um comprador que esteja disposto a desembolsar o valor pode adquiri-lo. Com um mercado cada vez mais competitivo, o desafio do vendedor passa a ser como precificar seus produtos.

% preços mais elevados aumentam a receita unitária, mas podem reduzir o conjunto de compradores dispostos a adquirir determinado produto; por outro lado, preços mais baixos ampliam o número de vendas possíveis, porém diminuem o ganho obtido em cada transação.

Neste contexto, surge um dilema econômico: se o vendedor aumenta o preço unitário, visando o lucro, pode diminuir as vendas; se diminuir demais o preço, visando aumentar as vendas, pode prejudicar o lucro. Assim, do ponto de vista do vendedor, o problema consiste em definir preços para um conjunto de itens de forma a equilibrar essas forças antagônicas e maximizar a receita total obtida.

% Com o objetivo de maximizar a receita do vendedor, define-se o \textit{problema da compra mínima} como segue. Seja $I$ um conjunto de itens, $B$ um conjunto de compradores e $V$ uma matriz em que cada elemento $v_{ib}$ representa o valor que o comprador $b \in B$ está disposto a pagar pelo item $i \in I$. O problema consiste em determinar um conjunto de preços $P$, em que cada elemento $p_i$ corresponde ao preço atribuído ao item $i$. Além disso, o comportamento do comprador é conhecido: ele comprará um único produto mais barato viável, caso contrário, ele não comprará nada.

O \textit{Problema da Compra Mínima} foi pouco explorado na literatura sob a perspectiva de algoritmos exatos. Além disso, os resultados experimentais indicam que há espaço para aprimoramentos nestes métodos, bem como para o desenvolvimento de heurísticas capazes de tratar, em tempo hábil, instâncias de grande porte do problema.

Portanto, este trabalho tem como objetivo propor melhorias às formulações matemáticas existentes, com o intuito de fortalecer os modelos de Programação Linear Inteira e Inteira Mista, bem como desenvolver heurísticas capazes de obter soluções de boa qualidade em tempo hábil para instâncias de grande porte do problema. 

% eficientes para a resolução de instâncias de grande porte do problema em tempo hábil.
\end{abstract}